\documentclass[fleqn]{ctexart}
% 数学支持（提供 \mathbb、\dfrac 等）
\usepackage{amsmath,amssymb}
% 页面与版式设置：更小的页边距与顶部空白
\usepackage[a4paper,top=15mm,bottom=20mm,left=12mm,right=12mm]{geometry}
% 自动换行与字偶距优化
\usepackage{microtype}
\microtypesetup{tracking=true,protrusion=true,final}
% 在需要时放宽断行以避免溢出（数值可按需调整）
\setlength{\emergencystretch}{2em}
% 列表控制，便于让(1)后直接跟内容
\usepackage{enumitem}
% 使章节标题左对齐（同时保留 ctex 默认样式）
\ctexset{section={format+=\raggedright}}
%1.3倍行距
\renewcommand{\baselinestretch}{1.3}
% 数学左对齐缩进为0
\setlength{\mathindent}{0pt}
% 增大矩阵的行距
\renewcommand{\arraystretch}{1.3}
\pagestyle{empty}
\begin{document}
\section*{题目1}
\begin{enumerate}[label=(\arabic*),leftmargin=0pt,itemindent=2em]
    \item A page fault（缺页）
    
    程序需要的数据不在内存里，必须去硬盘里取

    发生故障，属于同步异常
    \item You receive an E-mail from the Internet（互联网上收到一封邮件）
    
    这是一个独立于当前执行指令的外部事件
    
    发生中断，属于异步异常
    \item You’re running a command “kill -9 <pid>” in your shell（运行一句指令）
    
    使用了系统调用，有意触发该异常
    
    发生陷阱，属于同步异常
    \item The memory of your pc corrupted（计算机内存损坏）
    
    硬件损坏，无法通过处理程序进行修复
    
    发生终止，属于同步异常
\end{enumerate}

\section*{题目2}
\begin{verbatim}
main:
    movq $37, %rax
    movq $2024,%rdi
    movq $3, %rsi
    syscall

    movq $1, %rax
    movq $0, %rdi
    syscall
\end{verbatim}
把系统调用编号赋值给rax，其余参数按照顺序赋值给对应的寄存器，即可实现syscall。这里面的系统调用编号按照题目要求使用的是32位 Linux 系统的编号。
\end{document}